%=======================================================================
\section{Problem Statement --- Image Classification with Machine Learning}

\subsection{Clinical motivation}
Brain tumors are among the most life-threatening conditions, and the \textbf{type of tumor dictates
the treatment}. Gliomas, meningiomas, and pituitary tumors differ in origin, location, growth
pattern, and prognosis, so distinguishing them from an MRI scan is a decisive early step. In practice
this reading is done manually by radiologists, which is \textbf{time-consuming, requires scarce
expertise, and is subject to inter-observer variability}. An automated classifier can act as a fast,
consistent second opinion and a triage aid.

\begin{figure}[htbp]
  \centering
  \includegraphics[width=\textwidth]{dataset}
  \caption{Dataset overview for the classification task. \textbf{Left:} class distribution
    after an 80/20 stratified split (train in blue, test in orange) for the three tumor
    types. \textbf{Right:} mosaic of preprocessed training slices showing acquisition
    variability in orientation, contrast, and tumor appearance --- the visual diversity the
    classifier must handle.}
  \label{fig:dataset-overview}
\end{figure}

\subsection{Why this is a hard machine-learning problem}
Casting tumor identification as an image-classification task exposes several difficulties that the
rest of the report is designed to address:

\begin{itemize}[leftmargin=1.4em,itemsep=2pt]
  \item \textbf{High dimensionality.} A modest $128\times128$ grayscale image already has $16{,}384$
        raw pixels. Feeding raw pixels to a classical model invites the curse of dimensionality and
        overfitting.
  \item \textbf{Subtle, texture-based differences.} Tumor classes are not separated by simple colour
        or shape cues; the discriminative information lives in fine texture, intensity heterogeneity,
        and edge structure --- properties a single feature type rarely captures alone.
  \item \textbf{Acquisition variability.} MRI slices vary in contrast, orientation, field of view,
        and the amount of black background around the brain. Without normalisation, the model learns
        the scanner rather than the disease.
  \item \textbf{Limited, imbalanced data.} Medical datasets are small relative to natural-image
        datasets; class balance depends on the source corpus (Figure~\ref{fig:dataset-overview}).
        Stratification and macro-averaged metrics are therefore required.
  \item \textbf{The need for trust.} A clinical model must be not only accurate but interpretable and
        statistically credible --- a single accuracy number on one test split is not enough evidence.
\end{itemize}

\subsection{Our design choice: classical ML with handcrafted features}
Deep CNNs are the obvious modern tool, but they need large datasets, heavy compute, and behave as
black boxes. This project deliberately takes the \textbf{classical route}: engineer interpretable
descriptors that encode domain knowledge (texture, frequency content, gradient structure), then let
well-understood classifiers learn from a compact, meaningful feature space. This is fast to train,
runs on a CPU, is fully reproducible, and --- crucially for a clinical context --- every feature has
a clear physical meaning.

\begin{figure}[htbp]
  \centering
  \includegraphics[width=\textwidth]{problem_pipeline}
  \caption{Classical machine-learning pipeline adopted in this project: one preprocessed MRI
    slice is converted into compact handcrafted descriptors, fused and scaled, then classified
    by a tuned classical model. This design trades end-to-end deep learning for interpretability,
    low compute cost, and reproducibility.}
  \label{fig:problem-pipeline}
\end{figure}

\begin{notebox}[Task summary:]
Input $=$ one preprocessed MRI slice. Output $=$ one of \{Glioma, Meningioma, Pituitary tumor\}.
Method $=$ handcrafted multi-descriptor features $\rightarrow$ fusion $\rightarrow$ tuned classical
classifier, validated statistically.
\end{notebox}
